\clearpage\pagestyle{otMain}\setlength{\otBlocEspacement}{2.5mm}\fontsize{11}{15}\selectfont\corrigesChapitre{3}{Découvrir les nombres relatifs}
\begin{blocLivre}{cor-3-1}\begin{corrige}{1}{Lire une graduation}\demarche{Identifier l’unité et l’origine.}
Chaque intervalle vaut une unité. On lit $\res{A(-3),\ B(1),\ C(4)}$. Le signe dépend du côté de zéro.
\label{fig-cor-3-1}\figureGeo[corrige]{rel-1}
\end{corrige}\end{blocLivre}
\begin{blocLivre}{cor-3-2}\begin{corrige}{2}{Placer deux nombres}\demarche{Compter à partir de zéro.}
D se place deux unités à gauche de 0 ; E trois unités à droite. Les positions sont ajoutées en vert.
\label{fig-cor-3-2}\figureGeo[corrige]{rel-2}
\end{corrige}\end{blocLivre}
\begin{blocLivre}{cor-3-3}\begin{corrige}{3}{Trouver les opposés}\demarche{Conserver la distance à zéro.}
Ils sont $\res8$, $\res{-2{,}4}$ et $\res0$. Zéro est son propre opposé.
\end{corrige}\end{blocLivre}
\begin{blocLivre}{cor-3-4}\begin{corrige}{4}{Distances à zéro}\demarche{Lire une distance à zéro.}
Le premier point est à $\res{6{,}5}$ unités de zéro, le second à $\res4$ unités. Ce sont les valeurs absolues de $-6{,}5$ et de 4 : une distance est positive ou nulle.
\end{corrige}\end{blocLivre}
\begin{blocLivre}{cor-3-5}\begin{corrige}{5}{Un premier ordre}\demarche{Utiliser la droite orientée.}
Un négatif est inférieur à un positif : $\res{-5<2}$. Parmi les négatifs, $-9$ est plus à gauche : $\res{-9<-4}$.
\end{corrige}\end{blocLivre}
\begin{blocLivre}{cor-3-6}\begin{corrige}{6}{Ranger des décimaux}\demarche{Séparer les signes puis comparer.}
Les négatifs viennent d’abord ; $-1{,}2<-1{,}02$. Ordre croissant : $\res{-1{,}2<-1{,}02<0<0{,}8<1{,}2}$. Ordre décroissant : $\res{1{,}2>0{,}8>0>-1{,}02>-1{,}2}$.
\end{corrige}\end{blocLivre}
\begin{blocLivre}{cor-3-7}\begin{corrige}{7}{Encadrer un négatif}\demarche{Lire l’ordre, pas seulement les chiffres.}
Il est à gauche de $-3$ mais à droite de $-4$ : $\res{-4<-3{,}7<-3}$.
\end{corrige}\end{blocLivre}
\begin{blocLivre}{cor-3-8}\begin{corrige}{8}{Affiner une graduation}\demarche{Déterminer la valeur d’un pas.}
Chaque petit intervalle vaut $1/4=0{,}25$. Les points sont $\res{F(-0{,}75)}$ et $\res{G(1{,}5)}$.
\label{fig-cor-3-8}\figureGeo[corrige]{rel-8}
\end{corrige}\end{blocLivre}
\begin{blocLivre}{cor-3-9}\begin{corrige}{9}{Zéro et les signes}\demarche{Appliquer les définitions exactes.}
Les nombres positifs sont $\res{0\text{ et }5}$. Le seul strictement positif est $\res5$. Le mot « strictement » exclut zéro.
\end{corrige}\end{blocLivre}
\begin{blocLivre}{cor-3-10}\begin{corrige}{10}{Même distance}\demarche{Chercher des deux côtés de zéro.}
À quatre unités de zéro : $\res{-4\text{ et }4}$, deux nombres opposés. À distance nulle : $\res0$ seulement.
\end{corrige}\end{blocLivre}
\begin{blocLivre}{cor-3-11}\begin{corrige}{11}{Comparer des altitudes}\demarche{Identifier le sens du classement.}
L’ordre est $\res{-25<-8<0<12}$. Les deux premiers sont sous le niveau de la mer ; 0 m est au niveau de référence.
\end{corrige}\end{blocLivre}
\begin{blocLivre}{cor-3-12}\begin{corrige}{12}{Trouver un intermédiaire}\demarche{Ajouter un rang décimal.}
Par exemple $\res{-2{,}38\text{ et }-2{,}35}$. On vérifie $-2{,}4<-2{,}38<-2{,}35<-2{,}3$. Les bornes ne sont pas admises.
\end{corrige}\end{blocLivre}
\begin{blocLivre}{cor-3-13}\begin{corrige}{13}{Des contraintes entières}\demarche{Lire chaque borne.}
On exclut $-3$ et on inclut 2 : $\res{-2,-1,0,1,2}$. Les signes $<$ et $\le$ n’imposent pas la même chose.
\end{corrige}\end{blocLivre}
\begin{blocLivre}{cor-3-14}\begin{corrige}{14}{Une fausse règle}\demarche{Choisir deux signes différents.}
$-8$ est à 8 unités de zéro et 3 à seulement 3 unités, pourtant $\res{-8<3}$. La distance à zéro, c’est-à-dire la valeur absolue, ne suffit pas pour comparer : il faut aussi tenir compte du signe.
\end{corrige}\end{blocLivre}
\begin{blocLivre}{cor-3-15}\begin{corrige}{15}{Classement et opposés}\demarche{Réfléchir l’intervalle autour de zéro.}
Les positions possibles de $a$ sont entre $-6$ et $-2$. Leurs symétriques par rapport à zéro sont entre 2 et 6 : $\res{2<-a<6}$. Ici, $-a$ est donc $\res{\text{positif}}$ : l’écriture $-a$ ne désigne pas toujours un nombre négatif.
\label{fig-cor-3-15}\figureGeo[corrige]{rel-15}
\end{corrige}\end{blocLivre}
